\section{Coupled geometry of the binary measurement ball}\label{sec:coupled74}
We now encode the visibility as well as the direction. More generally,
we give a local geometric criterion for arbitrary regular subfamilies
of the unbiased binary measurement ball. The parameter space in this
criterion is the identifiable image of the family, not a redundant
choice of angular coordinates.

Let $\mathbb B=\{x\in\R^3:|x|\le1\}$. With ordered outcomes put
\begin{equation}\label{eq:ballfamily74}
 E_x^\pm=\tfrac12(I\pm x\cdot\sigma),\qquad
 \mathcal M_x^\pm(X)=\tr(E_x^\pm X),\qquad x\in\mathbb B.
\end{equation}
Here $\sigma=(\sigma_1,\sigma_2,\sigma_3)$ is the Pauli triple. We use
$d_N$ with its unhalved trace-norm convention and all common
reference-assisted adaptive testers with at most $N$ calls, including
feedback and bounded public stopping. The channel outputs the ordered
classical outcome and no residual quantum system. Define
\begin{equation}\label{eq:weight74}
 a_N(r)=\sqrt{\frac{N}{1-r^2+N^{-1}}},\qquad
 w_N(x)=a_N(|x|),\qquad N\ge1.
\end{equation}
The cutoff $N^{-1}$ is part of a finite-use estimate. It is not a
positive-noise assumption on the target.

\begin{theorem}[Finite-use metric on the measurement ball]\label{thm:ballmetric74}
For $x,y\in\mathbb B$, let $s=\min(|x|,|y|)$ and
$D_N(x,y)=a_N(s)|x-y|$. Then
\begin{equation}\label{eq:ballmetric74}
 \frac1{256}\min\{1,D_N(x,y)\}
 \le d_N(\mathcal M_x,\mathcal M_y)
 \le\min\{2,6D_N(x,y)\}.
\end{equation}
In particular $d_1(\mathcal M_x,\mathcal M_y)=|x-y|$. The constants
in \eqref{eq:ballmetric74} are independent of the visibility, direction,
horizon and pair separation. The lower bound is attained up to these
constants by pair-dependent nonadaptive tests; it does not assert the
existence of one estimator for the whole ball.
\end{theorem}
This supplies a finite-distance local comparison even at erasure and
at projective effects. The comparison is not obtained by integrating a
Fisher metric through a singular boundary. We first prove its radial
part, for which there is an exact classical reduction.

\begin{lemma}[Binomial distance with a support cutoff]\label{lem:radial74}
For $0\le s\le r\le1$, a unit $u$, and $P_t=\operatorname{Bern}((1-t)/2)$,
\begin{equation}\label{eq:radialequality74}
 d_N(\mathcal M_{ru},\mathcal M_{su})
       =\|P_r^{\otimes N}-P_s^{\otimes N}\|_1.
\end{equation}
Writing $R=a_N(s)(r-s)$, this distance lies between
$\frac1{64}\min\{1,R\}$ and $\min\{2,3R\}$.
\end{lemma}
\begin{proof}
First perform the projective measurement along $u$ and then flip its
output with probability $(1-t)/2$. For any fixed adaptive tester, all
its outputs are a common channel of $N$ independent flip bits. This
proves the upper inequality in \eqref{eq:radialequality74}. Repeated
input of the $+u$ eigenstate gives exactly these bits and proves equality.

Put $p=(1-r)/2$, $q=(1-s)/2$, and $\Delta=q-p$. The case $\Delta=0$
is immediate. Otherwise $0<q\le1/2$. Let
\[
 k=\min\{N,\lfloor(2q)^{-1}\rfloor\},\qquad m=\lfloor N/k\rfloor.
\]
Then $k\ge\frac14\min\{N,q^{-1}\}$ and $m\ge N/(2k)$.
In a block of $k$ trials the no-success event has probabilities
$e_p=(1-p)^k$ and $e_q=(1-q)^k$. Their difference obeys
\[
 e_p-e_q\ge k\Delta(1-q)^{k-1}\ge k\Delta/2.
\]
The second inequality follows from $(k-1)q\le1/2$ and Bernoulli's
inequality. A common randomized processing maps a Bernoulli parameter
$e$ to $b+ce$, where
\[
 c=\frac1{\max\{e_p+e_q,2-e_p-e_q\}}\ge\tfrac12,
 \qquad b=\tfrac12(1-c(e_p+e_q)).
\]
Indeed $b\ge0$, $b+c\le1$, and the processed pair has parameters
$(1\pm c(e_p-e_q))/2$. Apply
Lemma~\ref{lem:bernoullimajority73} to the $m$ independent processed
blocks. Contraction under processing gives
\[
 \|P_r^{\otimes N}-P_s^{\otimes N}\|_1
 \ge\frac1{16\sqrt2}\min\{1,\Delta\sqrt{N\min(N,q^{-1})}\}.
\]
Since $1-s^2=2q(1+s)$, we have
$a_N(s)\le\sqrt{N\min(N,q^{-1})}$. Also $\Delta=(r-s)/2$.
The asserted lower constant $1/64$ follows.

For the upper bound a one-bit coupling gives $N(r-s)$. Let
$F=\sqrt{pq}+\sqrt{(1-p)(1-q)}$. Directly,
\[
 1-F=\tfrac12\bigl[(\sqrt p-\sqrt q)^2+
                 (\sqrt{1-p}-\sqrt{1-q})^2\bigr]
 \le\frac{\Delta^2}{2q(1-q)}
 =\frac{(r-s)^2}{2(1-s^2)}.
\]
The product-fidelity inequality therefore gives the second bound
$2\sqrt N(r-s)/\sqrt{1-s^2}$ when $s<1$. For $z=N(1-s^2)\ge0$,
\[
 \min\{1,2/\sqrt z\}\sqrt{z+1}\le\sqrt5<3.
\]
Combining the two upper estimates proves the claim, without division
by zero at the identical endpoint $r=s=1$.
\end{proof}

\begin{proof}[Proof of Theorem~\ref{thm:ballmetric74}]
Assume $x=ru$, $y=sv$, $r\ge s>0$, and write $h\in[0,\pi]$ for the
angle between $u,v$. Any two directions lie on a great circle; common
unitary conjugation of all input systems preserves $d_N$. Thus the
fixed-visibility theorem \ref{thm:noisycoding73} applies with
$K=K_{N,s}$. Notice
\begin{equation}\label{eq:angularweight74}
 sa_N(s)\le K\le\sqrt2\,sa_N(s).
\end{equation}
Let $d=d_N(\mathcal M_{ru},\mathcal M_{sv})$ and
$d_{\rm r}=d_N(\mathcal M_{ru},\mathcal M_{su})$.
Inputting the $+u$ eigenstate gives Bernoulli parameters
$p=(1-r)/2$ and $q'=(1-s\cos h)/2\ge q=(1-s)/2$.
For fixed $p$, the total variation of binomial products increases
with the larger parameter. To see this without a differentiability
assumption, the binomial likelihood ratio is monotone in the success
count; its optimal event is a lower count interval. Each such event
has probability decreasing in the other parameter. Endpoint cases
follow by continuity. Consequently \eqref{eq:radialequality74} gives
$d_{\rm r}\le d$. The triangle inequality now yields
\[
 d_N(\mathcal M_{su},\mathcal M_{sv})\le d_{\rm r}+d\le2d.
\]
Set $A=a_N(s)(r-s)$ and $B=sa_N(s)h$. The radial lemma and the
fixed-visibility lower bound imply
\[
 d\ge\tfrac1{64}\min\{1,A\},\qquad
 d\ge\tfrac1{128}\min\{1,B\}.
\]
Hence $d\ge\frac1{256}\min\{1,A+B\}$. Since
$|ru-sv|\le(r-s)+sh$, this proves the required lower bound.
For the upper bound the radial and angular estimates give
$d\le3A+\sqrt2 B$. Moreover $r-s\le|ru-sv|$ and
$sh\le(\pi/2)|ru-sv|$, by
$|ru-sv|^2=(r-s)^2+4rs\sin^2(h/2)$. Thus
$d\le(3+\pi/\sqrt2)D_N<6D_N$.
When $s=0$, the radial lemma proves both bounds, independently of the
undefined direction at the origin. The one-use identity follows, as
in Lemma~\ref{lem:noisyupper73}, from the eigenvalues
$\pm|x-y|/2$ of $E_x^+-E_y^+$ and reference trace-norm duality.
\end{proof}

\subsection{A local volume criterion and boundary contact}
For compact $X\subset\mathbb B$ let $\mathcal C_N(X,\delta)$ allow
arbitrary legal memoryless centres of the binary qubit interface, and
let $\mathcal C_N^X(X,\delta)$ require centres in
$\{\mathcal M_x:x\in X\}$. The measure in the next theorem is ordinary
Euclidean local geometry on the identifiable set $X$.

\begin{theorem}[Weighted local covering criterion]\label{thm:weightedcover74}
Suppose a finite Borel measure $\mu$ on a nonempty compact
$X\subset\mathbb B$ satisfies, for some $k>0$, $r_0>0$ and $c_0,C_0>0$,
\begin{equation}\label{eq:ahlfors74}
 c_0 t^k\le\mu(X\cap B(x,t))\le C_0t^k
           \quad(x\in X,\ 0<t\le r_0).
\end{equation}
There are $c_X,C_X,\delta_X>0$, depending only on these regularity
constants, such that, for every $N\ge1$ and $0<\delta\le\delta_X$,
\begin{equation}\label{eq:weightedcover74}
 c_X\delta^{-k}\int_X w_N(x)^k\,d\mu(x)
 \le\mathcal C_N(X,\delta)
 \le\mathcal C_N^X(X,\delta)
 \le C_X\delta^{-k}\int_X w_N(x)^k\,d\mu(x).
\end{equation}
No map retracting arbitrary centres onto $X$ is assumed. In a
parametrized family the set $X$ is its image under \eqref{eq:ballfamily74};
indistinguishable parameter values are not charged twice.
\end{theorem}
\begin{proof}
Put $e_N(x)=1-|x|^2+N^{-1}$, so $e_N\ge N^{-1}$ and
$|e_N(x)-e_N(y)|\le2|x-y|$. If
$|x-y|\le t/w_N(x)$ and $t\le1/4$, then
\[
 |e_N(x)-e_N(y)|\le2t\sqrt{e_N(x)/N}\le2t e_N(x).
\]
Thus the weights are comparable on these small Euclidean balls.
Conversely, if $d_N(\mathcal M_x,\mathcal M_y)\le t<1/256$,
Theorem~\ref{thm:ballmetric74} implies
$|x-y|\le256t/a_N(s)$, $s=\min(|x|,|y|)$. Applying the same estimate
first at the point of smaller norm shows comparability to $w_N(x)$
when, for example, $t\le2^{-14}$. Consequently absolute $c,C>0$
satisfy the two local inclusions
\[
 X\cap B(x,ct/w_N(x))\ \subseteq\ B_{d_N}(x,t)\cap X
 \ \subseteq\ X\cap B(x,Ct/w_N(x)).
\]
Reduce $t$ by a fixed multiple of $r_0$ if necessary. Since
$w_N(x)\ge1/\sqrt2$, all Euclidean radii used here lie below $r_0$.
The weighted measure $d\nu_N=w_N^k d\mu$ of either operational ball
is bounded above and below by fixed multiples of $t^k$, by
\eqref{eq:ahlfors74} and weight comparability.

A maximal $\delta$-separated set in the compact metric space
$(X,d_N)$ covers $X$ by radius-$\delta$ balls. Its radius-$\delta/3$
balls are disjoint, proving the upper bound in
\eqref{eq:weightedcover74}. A maximal $3\delta$-separated set covers
$X$ by radius-$3\delta$ balls, so it has at least
$c\nu_N(X)\delta^{-k}$ points. Any radius-$\delta$ ball about any
legal centre contains at most one of these points, by the triangle
inequality. This gives the lower bound for unrestricted centres.
Compactness and the one-use identity ensure that these are genuine
metrics on $X$; the finite hybrid bound gives the required continuity.
\end{proof}

\begin{corollary}[Boundary-contact law]\label{cor:contact74}
In Theorem~\ref{thm:weightedcover74}, put
$F(t)=\mu\{x\in X:1-|x|^2\le t\}$. Suppose $F(0)=0$ and
$c t^\alpha\le F(t)\le C t^\alpha$ for $0<t\le t_0$, with
$\alpha>0$. Uniformly on the same small-error range,
\begin{equation}\label{eq:contact74}
 \mathcal C_N(X,\delta)\asymp_X\delta^{-k}
 \begin{cases}
 N^{k/2},&\alpha>k/2,\\
 N^{k/2}\log(N+2),&\alpha=k/2,\\
 N^{k-\alpha},&0<\alpha<k/2.
 \end{cases}
\end{equation}
If instead $F(0)>0$, the law is $\asymp_X\delta^{-k}N^k$.
\end{corollary}
\begin{proof}
Write $b=k/2$ and $\tau=N^{-1}$. Stieltjes integration by parts gives
\[
 \int_{[0,t_0]}(t+\tau)^{-b}\,dF(t)
 =(t_0+\tau)^{-b}F(t_0)
       +b\int_0^{t_0}\frac{F(t)}{(t+\tau)^{b+1}}\,dt
\]
when $F(0)=0$. Substitute the two bounds on $F$ and split the last
integral at $\tau$. It is bounded, logarithmic, or comparable to
$\tau^{\alpha-b}$ according as $\alpha>b$, $\alpha=b$, or
$\alpha<b$. The part beyond $t_0$ is bounded, and multiplication by
$N^b$ proves \eqref{eq:contact74}. Bounded $N$ is absorbed into the
constants. An atom at depth zero contributes exactly $N^kF(0)$ to
the weighted volume; $w_N\le N$ supplies the matching upper bound.
\end{proof}
For example, in a regular $k$-dimensional chart with a boundary subset
of codimension $\ell$, a two-sided relation between the support deficit
$1-|x|^2$ and the $m$th power of transverse distance yields
$\alpha=\ell/m$, provided the transverse volume has order $t^\ell$.
This is a geometric condition to check, not an automatic regularity
assertion. It distinguishes ordinary interior accumulation, critical
logarithmic enhancement and boundary-dominated coherent accumulation.

\begin{corollary}[Joint visibility and direction]\label{cor:jointball74}
Let $\mathbb B_d$ be the closed unit ball in a fixed $d$-dimensional
linear subspace of $\R^3$, $d=1,2,3$. For an absolute small-error range,
\begin{equation}\label{eq:jointball74}
 \mathcal C_N(\mathbb B_d,\delta)\asymp
 \begin{cases}
 \sqrt N\,\delta^{-1},&d=1,\\
 N\log(N+2)\,\delta^{-2},&d=2,\\
 N^2\,\delta^{-3},&d=3.
 \end{cases}
\end{equation}
For the coordinate disk and full ball, rational target data admit exact
rational family-centred codes attaining these orders. In particular,
the optimal fixed-length payloads for $d=2,3$ are respectively
\[
 \log_2N+\log_2\log(N+2)+2\log_2(1/\delta)+O(1),
 \qquad 2\log_2N+3\log_2(1/\delta)+O(1).
\]
Visibility is a charged target coordinate in these statements.
\end{corollary}
\begin{proof}
Lebesgue measure on $\mathbb B_d$ satisfies \eqref{eq:ahlfors74} with
$k=d$, and its boundary-depth distribution has exponent $\alpha=1$.
Apply Corollary~\ref{cor:contact74}. The exact rational upper
construction follows from Theorem~\ref{thm:jointcodec74} below. In a
coordinate line the direct scalar grid in the same construction gives
the stated one-dimensional order.
\end{proof}

\subsection{Exact joint coding without irrational normalization}
We give a rational construction for $d=2,3$ that takes the Cartesian
Bloch vector as input. Its norm need not be rational. Fix public
$N\ge1$ and rational $0<\delta\le1/4$. Set
\begin{equation}\label{eq:jointgrids74}
 B=\left\lceil\sqrt{64N/\delta^2}\right\rceil,\quad
 t_i=i/B,\quad r_i=\frac{2t_i}{1+t_i^2},\quad
 A_i=\max\left\{1,\left\lceil\sqrt{64(d-1)K_{N,r_i}^2/\delta^2}
                                      \right\rceil\right\}\quad(i>0).
\end{equation}
The radial layer $i=0$ has one codeword. At $i>0$ use $2d$ signed
stereographic charts and the grid $\{-A_i,\ldots,A_i\}^{d-1}/A_i$.
For chart axis $j$ and sign $s$, the map is
\begin{equation}\label{eq:spherechart74}
 u_j=s\frac{1-|z|^2}{1+|z|^2},\qquad
 u_\ell=\frac{2z_\ell}{1+|z|^2}\quad(\ell\ne j).
\end{equation}
Every decoded centre is $r_i u\in\mathbb B_d$ with rational entries.
The exact capacity, including possibly redundant chart words, is
\begin{equation}\label{eq:jointcapacity74}
 L_{N,\delta,d}=1+\sum_{i=1}^{B}2d(2A_i+1)^{d-1}.
\end{equation}
A single fixed-length integer index in this union of layers charges
the visibility digit as well as the direction digits.

\begin{theorem}[Rational joint codec]\label{thm:jointcodec74}
For rational $x\in\mathbb B_d$, $d=2,3$, an exact rational algorithm
returns an index in \eqref{eq:jointcapacity74}. Its decoded measurement
is within $\delta$ in $d_N$, and every index decodes to a legal rational
measurement. There are absolute constants such that
\[
 L_{N,\delta,2}\le C N\log(N+2)\delta^{-2},\qquad
 L_{N,\delta,3}\le C N^2\delta^{-3}.
\]
The stored payload is $\lceil\log_2L_{N,\delta,d}\rceil$ bits. The
JSON envelope, encoder workspace, public $d,N,\delta$ and expanded
matrices are not this payload. No input visibility is public, and no
floating-point operation decides a chart or grid digit.
\end{theorem}
\begin{proof}
Write $r=|x|$, $\beta=\arcsin r$, $t=\tan(\beta/2)\in[0,1]$.
Choose a nearest $i/B$ to $t$, resolving ties towards zero. Then
$|\beta-2\arctan(i/B)|\le1/B$. The radial flip programmes have root
fidelity $\cos((\beta-\beta_i)/2)$, and hence their $N$-copy distance
is at most $\sqrt N|\beta-\beta_i|\le\delta/8$.
When $i=0$, decode the unique erased measurement and stop.
Otherwise choose a maximal-absolute-coordinate axis of $x$, with its
sign. The exact chart parameters of $x/r$ are
\[
 z_\ell=\frac{x_\ell}{r+|x_j|},\qquad \ell\ne j.
\]
They lie in $[-1,1]$. Round each to a nearest grid point at denominator
$A_i$. The differential norm of \eqref{eq:spherechart74}, measured
in spherical length, is $2/(1+|z|^2)\le2$. Thus the angular error is
at most $\sqrt{d-1}/A_i$ and its operational cost is at most $\delta/8$,
by the fixed-visibility upper bound and \eqref{eq:jointgrids74}.
The triangle inequality leaves at least three quarters of the
requested budget unused.

All comparisons above can be made in the rationals. If $Q=\sum x_j^2$,
comparison of $t$ with a rational $v\in[0,1]$ is comparison of $Q$
with $[2v/(1+v^2)]^2$. For a direction digit compare
$b/(\sqrt Q+a)$ with rational $v\ge0$, where $a=|x_j|$ and
$b=|x_\ell|$. First test the sign of $b-va$. If it is negative, the
comparison is negative; otherwise compare $(b-va)^2$ and $v^2Q$.
The sign test is essential before squaring. Binary integer searches
therefore determine every nearest grid digit, including ties, without
computing an irrational normalization. Integer ceiling square roots
suffice for both grids.

It remains to bound the union capacity. Here $B\asymp\sqrt N/\delta$.
For $i\le B/2$, $K_{N,r_i}\le C\sqrt N$, which contributes at most
$C N^{d/2}\delta^{-d}$. For $i>B/2$, put $j=B-i$. Since
$1-r_i^2=(1-t_i^2)^2/(1+t_i^2)^2$, we have
\[
 K_{N,r_i}\le C\min\{N,\sqrt N B/(j+1)\},
\]
with an adjusted absolute constant also at $j=0$.
Split the sum at $j+1=B/\sqrt N\asymp\delta^{-1}$.
For $d=2$, the remaining sum is harmonic and contributes at most
$C N\delta^{-2}\log(N+2)$. For $d=3$, the sum of $(j+1)^{-2}$
contributes at most $C N^2\delta^{-3}$. The near-boundary layers have
these same or smaller orders. The added constants in $A_i$ and the
single erased word are absorbed into these bounds. This proves the
capacity and payload assertions. Counting layer sizes and ranking or
unranking the integer grid are finite exact algorithms, with cost
polynomial in the numerical radial grid and operand lengths; no
polynomial bound in the encoded length of $N$ or $1/\delta$ is asserted.
\end{proof}

\begin{remark}[Priority and interface]\label{rem:jointscope74}
Sedl\'ak--Ziman \cite[Section VI, Example 4, equation (37)]{SZ74}
already treat distinct white-noise levels in single-shot qubit
measurement discrimination. Pucha\l a--Pawela--Krawiec--Kukulski
\cite[Theorem 1]{PPKK74} identify the single-shot diamond distance
of projective measurements with an optimized unitary-channel distance.
Neither result is claimed anew here. Our conclusion is the uniform
finite-use metric comparison \eqref{eq:ballmetric74}, its weighted
covering criterion, and the resulting contact and exact rational
joint-description laws. The classical programme and entangled-block
mechanisms retain their earlier attribution. The geometric criterion
covers regular identifiable subsets of \eqref{eq:ballfamily74}, not
arbitrary biased POVMs or general disturbing instruments. The supplied
vector is encoded; no unknown measurement is learned. Pairwise testing,
a common estimator, a global environment seizer and a physical classical
simulation of unknown quantum inputs are different tasks.
\end{remark}
